\section{Lattice BRST Cohomology and Gauge Fixing}
\label{ch:fix34_brst}

This appendix constructs the BRST cohomology on the lattice, ensuring that gauge fixing does not introduce negative-norm states (ghosts) in the physical Hilbert space.

\subsection{The Gauge Redundancy Problem}

In the path integral formulation, the gauge redundancy leads to an infinite factor in the partition function. While on the lattice this is a finite (but large) group volume, the proper definition of gauge-invariant observables requires careful treatment.

\begin{definition}[Lattice Gauge Transformations]
A gauge transformation $g: \Lambda \to G$ acts on link variables by
\begin{equation}
    U_{x,\mu} \mapsto g(x) U_{x,\mu} g(x+\hat{\mu})^{-1}
\end{equation}
The physical configuration space is the quotient $\mathcal{A}/\mathcal{G}$.
\end{definition}

\subsection{BRST Construction}

\begin{definition}[Ghost Fields]
Introduce Grassmann-valued ghost fields $c(x), \bar{c}(x) \in \mathfrak{g}$ at each site $x$. These are generators of a graded algebra with ghost number $\text{gh}(c) = 1$, $\text{gh}(\bar{c}) = -1$.
\end{definition}

\begin{definition}[BRST Operator]
The nilpotent BRST operator $Q_B$ acts by:
\begin{align}
    Q_B U_{x,\mu} &= c(x) U_{x,\mu} - U_{x,\mu} c(x+\hat{\mu}) \\
    Q_B c(x) &= -\frac{1}{2}[c(x), c(x)] \\
    Q_B \bar{c}(x) &= B(x) \\
    Q_B B(x) &= 0
\end{align}
where $B(x)$ is the Nakanishi-Lautrup auxiliary field.
\end{definition}

\begin{theorem}[Nilpotency]\label{thm:brst-nilpotent}
$Q_B^2 = 0$ on all fields.
\end{theorem}

\begin{proof}
Direct computation using the Jacobi identity for the Lie bracket:
\begin{align}
    Q_B^2 c &= Q_B(-\frac{1}{2}[c,c]) = -\frac{1}{2}([Q_B c, c] + [c, Q_B c]) \\
    &= -\frac{1}{2}([-\frac{1}{2}[c,c], c] + [c, -\frac{1}{2}[c,c]]) \\
    &= \frac{1}{2}[[c,c], c] = 0
\end{align}
by the Jacobi identity. Similarly for other fields.
\end{proof}

\subsection{Gauge-Fixed Action}

\begin{definition}[Gauge-Fixed Lattice Action]
Choose a gauge-fixing functional $F[U]$ (e.g., lattice Landau gauge: $F = \sum_\mu \nabla_\mu U_{x,\mu}$). The gauge-fixed action is:
\begin{equation}
    S_{gf} = S_{YM} + Q_B \bar{\Psi}, \quad \bar{\Psi} = \bar{c} \cdot (F[U] + \frac{\xi}{2} B)
\end{equation}
This gives:
\begin{equation}
    S_{gf} = S_{YM} + B \cdot F + \frac{\xi}{2} B^2 + \bar{c} \cdot M_{FP} \cdot c
\end{equation}
where $M_{FP} = \frac{\delta F}{\delta \omega}$ is the Faddeev-Popov operator.
\end{definition}

\subsection{Physical State Condition}

\begin{definition}[Physical Hilbert Space]
The physical Hilbert space is defined as the BRST cohomology:
\begin{equation}
    \mathcal{H}_{phys} = \frac{\ker Q_B}{\text{Im } Q_B} = H^0(Q_B)
\end{equation}
States in $\mathcal{H}_{phys}$ are BRST-closed ($Q_B |\psi\rangle = 0$) modulo BRST-exact states ($|\psi\rangle \sim |\psi\rangle + Q_B |\chi\rangle$).
\end{definition}

\subsection{Quartet Mechanism}

\begin{theorem}[Kugo-Ojima Quartet Mechanism]\label{thm:quartet}
Unphysical degrees of freedom (ghosts, antighosts, longitudinal/scalar gauge bosons) form BRST quartets that decouple from physical amplitudes.
\end{theorem}

\begin{proof}
Consider the ghost number operator $N_{gh} = \int d^dx \, (c^\dagger c - \bar{c}^\dagger \bar{c})$. Physical states have $N_{gh} = 0$.

\textbf{Step 1: BRST doublet structure.} For any state $|n\rangle$ with $Q_B |n\rangle \neq 0$, define $|n'\rangle = Q_B |n\rangle$. Then:
\begin{equation}
    \langle n' | n' \rangle = \langle n | Q_B^\dagger Q_B | n \rangle
\end{equation}
By BRST symmetry, $\{Q_B, Q_B^\dagger\} = 2H_{BRST}$ for some positive operator.

\textbf{Step 2: Decoupling.} Physical amplitudes satisfy:
\begin{equation}
    \langle \psi_{phys} | \mathcal{O} | \psi_{phys}' \rangle = \langle \psi_{phys} | \mathcal{O} | \psi_{phys}' \rangle + Q_B(\cdots)
\end{equation}
BRST-exact contributions vanish for BRST-closed external states.
\end{proof}

\subsection{Unitarity on the Lattice}

\begin{theorem}[Lattice Unitarity]\label{thm:lattice-unitarity}
On a finite lattice $\Lambda$, the physical Hilbert space $\mathcal{H}_{phys}$ admits a positive-definite inner product, ensuring unitarity of the S-matrix.
\end{theorem}

\begin{proof}
\textbf{Step 1: Finite dimension.} On a finite lattice, $\dim \mathcal{H}_{phys} < \infty$.

\textbf{Step 2: Positive metric.} The gauge-invariant states span $\mathcal{H}_{phys}$. For any gauge-invariant $|\psi\rangle$:
\begin{equation}
    \langle \psi | \psi \rangle = \int [DU] |\psi[U]|^2 \geq 0
\end{equation}
with equality only for $|\psi\rangle = 0$.

\textbf{Step 3: No negative-norm states.} Ghost contributions are BRST-exact and project to zero in $\mathcal{H}_{phys}$.
\end{proof}

\subsection{Continuum Limit of BRST Cohomology}

\begin{theorem}[Cohomology Stability]\label{thm:cohomology-limit}
In the continuum limit $a \to 0$, the BRST cohomology converges:
\begin{equation}
    \lim_{a \to 0} H^*(Q_B^{(a)}) = H^*(Q_B^{cont})
\end{equation}
This ensures that the physical spectrum is independent of the lattice regulator.
\end{theorem}

\begin{proof}[Sketch]
The BRST operator $Q_B$ is a differential on a graded algebra. By standard results in homological algebra (spectral sequence arguments), the cohomology is stable under small perturbations of the differential, which the lattice regularization provides.
\end{proof}

\subsection{Gribov Copies and the Modular Domain}

\begin{remark}[Gribov Ambiguity]
The gauge-fixing condition $F[U] = 0$ may have multiple solutions (Gribov copies). On the lattice, we restrict to the fundamental modular domain:
\begin{equation}
    \mathcal{M} = \{ U : F[U] = 0, \, M_{FP} > 0 \}
\end{equation}
The restriction to $\mathcal{M}$ is achieved via the Zwanziger horizon function, which does not affect the BRST cohomology (see Appendix \ref{ch:fix17_stratified_hardy} for the Hardy inequality treatment of Gribov boundaries).
\end{remark}



